\documentclass[11pt]{article}
\newcommand{\ListaTitulo}{Lista 06 --- Pressupostos e Transformações}
% Preâmbulo compartilhado — MATD48, Listas de Exercícios 2026
% Usado por Lista{NN}.tex e Gabarito{NN}.tex via % Preâmbulo compartilhado — MATD48, Listas de Exercícios 2026
% Usado por Lista{NN}.tex e Gabarito{NN}.tex via % Preâmbulo compartilhado — MATD48, Listas de Exercícios 2026
% Usado por Lista{NN}.tex e Gabarito{NN}.tex via \input{preamble}

\usepackage[utf8]{inputenc}
\usepackage[T1]{fontenc}
\usepackage[brazil]{babel}
\usepackage{fullpage}
\usepackage{amsmath,amssymb}
\usepackage{enumitem}
\usepackage{xcolor}
\usepackage{fancyhdr}
\usepackage{titlesec}
\usepackage{listings}
\usepackage{hyperref}
\usepackage{booktabs}
\usepackage{graphicx}

\definecolor{matd48azul}{HTML}{0B4F6C}
\definecolor{matd48verde}{HTML}{2E8B57}

\titleformat{\section}{\normalfont\Large\bfseries\color{matd48azul}}{\thesection}{1em}{}
\titleformat{\subsection}{\normalfont\large\bfseries\color{matd48azul}}{\thesubsection}{1em}{}

\providecommand{\ListaTitulo}{Lista de Exercícios}

\setlength{\headheight}{14pt}
\pagestyle{fancy}
\fancyhf{}
\lhead{\small MATD48 --- Planejamento de Experimentos A}
\rhead{\small \ListaTitulo}
\cfoot{\thepage}
\renewcommand{\headrulewidth}{0.4pt}

\lstset{
  basicstyle=\ttfamily\small,
  breaklines=true,
  frame=single,
  columns=fullflexible,
  backgroundcolor=\color{gray!8},
  commentstyle=\color{matd48verde},
  keywordstyle=\color{matd48azul}\bfseries,
  language=R,
  extendedchars=true,
  literate=%
    {á}{{\'a}}1 {à}{{\`a}}1 {â}{{\^a}}1 {ã}{{\~a}}1
    {é}{{\'e}}1 {ê}{{\^e}}1 {í}{{\'i}}1 {ó}{{\'o}}1
    {ô}{{\^o}}1 {õ}{{\~o}}1 {ú}{{\'u}}1 {ç}{{\c c}}1
    {Á}{{\'A}}1 {À}{{\`A}}1 {Â}{{\^A}}1 {Ã}{{\~A}}1
    {É}{{\'E}}1 {Ê}{{\^E}}1 {Í}{{\'I}}1 {Ó}{{\'O}}1
    {Ô}{{\^O}}1 {Õ}{{\~O}}1 {Ú}{{\'U}}1 {Ç}{{\c C}}1
}

\hypersetup{colorlinks=true, linkcolor=matd48azul, urlcolor=matd48azul, citecolor=matd48azul}

\newcommand{\cabecalho}[2]{%
  \begin{center}
    {\Large\bfseries\color{matd48azul} #1}\\[0.3em]
    {\large #2}\\[0.3em]
    \rule{\linewidth}{0.6pt}
  \end{center}
  \vspace{0.5em}
}

\newenvironment{questao}[1]{\vspace{1em}\noindent\textbf{Questão #1.}\hspace{0.4em}}{\par}
\newenvironment{solucao}{\vspace{0.3em}\noindent\textit{Solução.}\hspace{0.4em}\color{black!85}}{\par\vspace{0.5em}}


\usepackage[utf8]{inputenc}
\usepackage[T1]{fontenc}
\usepackage[brazil]{babel}
\usepackage{fullpage}
\usepackage{amsmath,amssymb}
\usepackage{enumitem}
\usepackage{xcolor}
\usepackage{fancyhdr}
\usepackage{titlesec}
\usepackage{listings}
\usepackage{hyperref}
\usepackage{booktabs}
\usepackage{graphicx}

\definecolor{matd48azul}{HTML}{0B4F6C}
\definecolor{matd48verde}{HTML}{2E8B57}

\titleformat{\section}{\normalfont\Large\bfseries\color{matd48azul}}{\thesection}{1em}{}
\titleformat{\subsection}{\normalfont\large\bfseries\color{matd48azul}}{\thesubsection}{1em}{}

\providecommand{\ListaTitulo}{Lista de Exercícios}

\setlength{\headheight}{14pt}
\pagestyle{fancy}
\fancyhf{}
\lhead{\small MATD48 --- Planejamento de Experimentos A}
\rhead{\small \ListaTitulo}
\cfoot{\thepage}
\renewcommand{\headrulewidth}{0.4pt}

\lstset{
  basicstyle=\ttfamily\small,
  breaklines=true,
  frame=single,
  columns=fullflexible,
  backgroundcolor=\color{gray!8},
  commentstyle=\color{matd48verde},
  keywordstyle=\color{matd48azul}\bfseries,
  language=R,
  extendedchars=true,
  literate=%
    {á}{{\'a}}1 {à}{{\`a}}1 {â}{{\^a}}1 {ã}{{\~a}}1
    {é}{{\'e}}1 {ê}{{\^e}}1 {í}{{\'i}}1 {ó}{{\'o}}1
    {ô}{{\^o}}1 {õ}{{\~o}}1 {ú}{{\'u}}1 {ç}{{\c c}}1
    {Á}{{\'A}}1 {À}{{\`A}}1 {Â}{{\^A}}1 {Ã}{{\~A}}1
    {É}{{\'E}}1 {Ê}{{\^E}}1 {Í}{{\'I}}1 {Ó}{{\'O}}1
    {Ô}{{\^O}}1 {Õ}{{\~O}}1 {Ú}{{\'U}}1 {Ç}{{\c C}}1
}

\hypersetup{colorlinks=true, linkcolor=matd48azul, urlcolor=matd48azul, citecolor=matd48azul}

\newcommand{\cabecalho}[2]{%
  \begin{center}
    {\Large\bfseries\color{matd48azul} #1}\\[0.3em]
    {\large #2}\\[0.3em]
    \rule{\linewidth}{0.6pt}
  \end{center}
  \vspace{0.5em}
}

\newenvironment{questao}[1]{\vspace{1em}\noindent\textbf{Questão #1.}\hspace{0.4em}}{\par}
\newenvironment{solucao}{\vspace{0.3em}\noindent\textit{Solução.}\hspace{0.4em}\color{black!85}}{\par\vspace{0.5em}}


\usepackage[utf8]{inputenc}
\usepackage[T1]{fontenc}
\usepackage[brazil]{babel}
\usepackage{fullpage}
\usepackage{amsmath,amssymb}
\usepackage{enumitem}
\usepackage{xcolor}
\usepackage{fancyhdr}
\usepackage{titlesec}
\usepackage{listings}
\usepackage{hyperref}
\usepackage{booktabs}
\usepackage{graphicx}

\definecolor{matd48azul}{HTML}{0B4F6C}
\definecolor{matd48verde}{HTML}{2E8B57}

\titleformat{\section}{\normalfont\Large\bfseries\color{matd48azul}}{\thesection}{1em}{}
\titleformat{\subsection}{\normalfont\large\bfseries\color{matd48azul}}{\thesubsection}{1em}{}

\providecommand{\ListaTitulo}{Lista de Exercícios}

\setlength{\headheight}{14pt}
\pagestyle{fancy}
\fancyhf{}
\lhead{\small MATD48 --- Planejamento de Experimentos A}
\rhead{\small \ListaTitulo}
\cfoot{\thepage}
\renewcommand{\headrulewidth}{0.4pt}

\lstset{
  basicstyle=\ttfamily\small,
  breaklines=true,
  frame=single,
  columns=fullflexible,
  backgroundcolor=\color{gray!8},
  commentstyle=\color{matd48verde},
  keywordstyle=\color{matd48azul}\bfseries,
  language=R,
  extendedchars=true,
  literate=%
    {á}{{\'a}}1 {à}{{\`a}}1 {â}{{\^a}}1 {ã}{{\~a}}1
    {é}{{\'e}}1 {ê}{{\^e}}1 {í}{{\'i}}1 {ó}{{\'o}}1
    {ô}{{\^o}}1 {õ}{{\~o}}1 {ú}{{\'u}}1 {ç}{{\c c}}1
    {Á}{{\'A}}1 {À}{{\`A}}1 {Â}{{\^A}}1 {Ã}{{\~A}}1
    {É}{{\'E}}1 {Ê}{{\^E}}1 {Í}{{\'I}}1 {Ó}{{\'O}}1
    {Ô}{{\^O}}1 {Õ}{{\~O}}1 {Ú}{{\'U}}1 {Ç}{{\c C}}1
}

\hypersetup{colorlinks=true, linkcolor=matd48azul, urlcolor=matd48azul, citecolor=matd48azul}

\newcommand{\cabecalho}[2]{%
  \begin{center}
    {\Large\bfseries\color{matd48azul} #1}\\[0.3em]
    {\large #2}\\[0.3em]
    \rule{\linewidth}{0.6pt}
  \end{center}
  \vspace{0.5em}
}

\newenvironment{questao}[1]{\vspace{1em}\noindent\textbf{Questão #1.}\hspace{0.4em}}{\par}
\newenvironment{solucao}{\vspace{0.3em}\noindent\textit{Solução.}\hspace{0.4em}\color{black!85}}{\par\vspace{0.5em}}


\begin{document}

\cabecalho{Lista de Exercícios 06}{Verificação de pressupostos e transformações no DCA (Capítulo 3 / Aula 06)}

\noindent \textit{Entrega: até a próxima aula. Justifique todas as respostas --- nestas listas, a
justificativa vale mais que a resposta final. O gabarito completo está em \texttt{Gabarito06.pdf}.}

\begin{questao}{1} \textbf{(Psicologia --- interpretando um Q-Q plot)}

O Q-Q plot dos resíduos de uma ANOVA (tempo de reação sob três condições de distração)
acompanha bem a reta de referência na região central, mas na cauda superior direita alguns
poucos pontos se destacam claramente \emph{acima} da reta (resíduos maiores do que a normal
prevê), enquanto a cauda inferior esquerda segue a reta normalmente.

\begin{enumerate}[label=(\alph*)]
  \item Que tipo de desvio de normalidade esse padrão sugere? Relacione com o formato típico de
    dados de tempo de reação.
  \item O pesquisador tem $N=45$ observações. Um teste de Shapiro-Wilk aplicado a esses resíduos
    não rejeita a normalidade ($p=0{,}09$). Isso contradiz o que o Q-Q plot sugere? Como você
    reconciliaria as duas informações antes de decidir se a ANOVA é confiável?
  \item Supondo que o padrão do Q-Q plot reflita um problema real, qual transformação da família
    de potências você tentaria primeiro, e por quê?
\end{enumerate}
\end{questao}

\begin{questao}{2} \textbf{(Psicologia --- homogeneidade de variância)}

Um estudo compara o tempo de reação (ms) sob quatro níveis de privação de sono (0h, 12h, 24h,
36h), com $n=20$ participantes por grupo ($N=80$). Os desvios-padrão amostrais observados são,
respectivamente, $30$, $34$, $68$ e $81$ ms. O teste de Levene aplicado a esses dados fornece
$F_{3,76} = 5{,}87$, $p = 0{,}0012$.

\begin{enumerate}[label=(\alph*)]
  \item O teste de Levene rejeita a hipótese de homocedasticidade a $\alpha=0{,}05$? O que isso
    implica para a validade dos ICs e do teste $F$ da ANOVA, tais como definidos na Aula 04?
  \item Calcule a razão entre a maior e a menor variância amostral ($s^2_{36h}/s^2_{0h}$). Esse
    padrão --- variância crescendo junto com a severidade da privação de sono --- é consistente
    com erro aditivo ou multiplicativo? Justifique.
  \item Proponha uma transformação para estabilizar a variância e explique, em termos do padrão
    observado (variância cresce com a média), por que ela deveria funcionar.
\end{enumerate}
\end{questao}

\begin{questao}{3} \textbf{(Dedução --- por que o log estabiliza a variância sob erro multiplicativo)}

Suponha que, em vez do modelo aditivo usual, o tempo de reação siga um modelo com erro
\textbf{multiplicativo}: $Y_{ij} = \mu_i \cdot \exp(\varepsilon_{ij})$, com $\varepsilon_{ij}
\overset{iid}{\sim} N(0,\sigma^2)$ e $\mu_i > 0$ a "escala" típica do grupo $i$. Use o fato de
que, se $\varepsilon \sim N(0,\sigma^2)$, então $W=\exp(\varepsilon)$ (lognormal) tem $E[W] =
e^{\sigma^2/2}$ e $\mathrm{Var}(W) = \big(e^{\sigma^2}-1\big)e^{\sigma^2}$.

\begin{enumerate}[label=(\alph*)]
  \item Mostre que $\log Y_{ij} = \log\mu_i + \varepsilon_{ij}$ segue exatamente o modelo aditivo
    usual do DCA (Aula 04), com erro $N(0,\sigma^2)$ de variância \textbf{constante}, não
    dependendo de $i$.
  \item Mostre que, na escala original, $\mathrm{Var}(Y_{ij}) = \mu_i^2 \cdot C$, em que $C =
    \big(e^{\sigma^2}-1\big)e^{\sigma^2}$ é uma constante que não depende de $i$ --- ou seja, a
    variância na escala original é proporcional ao \textbf{quadrado} da média do grupo.
  \item Combine os itens (a) e (b) para explicar, em uma frase, por que grupos com médias maiores
    aparecem com variância amostral maior na escala original (o padrão da Questão 2), e por que a
    transformação log elimina esse padrão.
\end{enumerate}
\end{questao}

\begin{questao}{4} \textbf{(Ciência de dados --- interpretando Box-Cox)}

O procedimento de Box-Cox aplicado aos resíduos de um modelo ANOVA (tempo de carregamento de uma
página, três configurações de servidor) estima $\hat\lambda = -0{,}4$, com intervalo de
verossimilhança de 95\% igual a $(-1{,}1,\ 0{,}3)$.

\begin{enumerate}[label=(\alph*)]
  \item O intervalo inclui $\lambda=0$ (log)? E $\lambda=1$ (sem transformar)? O que cada resposta
    implica?
  \item Com base no item (a), qual transformação você recomendaria usar na prática, e por quê a
    interpretabilidade pesa nessa escolha (não só o valor que maximiza a verossimilhança)?
  \item Se, em vez disso, o intervalo de 95\% fosse $(-2{,}5,\ -1{,}8)$ --- não contendo nenhum
    valor "redondo" usual ($0$, $1/2$, $-1$, $1$) --- que estratégias você consideraria?
\end{enumerate}
\end{questao}

\begin{questao}{5} \textbf{(Integradora --- decidindo o que fazer)}

Uma pesquisadora mede o número de erros cometidos em uma tarefa de revisão de texto sob três
condições de iluminação (fraca, média, forte), com 15 participantes por condição. Os dados são
contagens (0, 1, 2, 3, ...), a maioria concentrada em valores baixos, com alguns participantes
cometendo bem mais erros que a média em cada grupo.

\begin{enumerate}[label=(\alph*)]
  \item Que diagnósticos você aplicaria, e em que ordem, antes de confiar em uma ANOVA padrão
    nesses dados?
  \item Dado que são dados de contagem, qual transformação da Aula 06 é a candidata natural, e
    por quê?
  \item Suponha que, mesmo após tentar as transformações usuais, o Q-Q plot dos resíduos ainda
    mostre desvios severos de normalidade. Sem entrar em detalhes técnicos (assunto da próxima
    seção do curso), que tipo de alternativa a pesquisadora deveria considerar, e o que essa
    alternativa muda conceitualmente em relação a comparar médias?
\end{enumerate}
\end{questao}

\end{document}
